\documentclass[tikz]{standalone}
\usepackage{../sty/phylogenetic-diagrams} % Include the style file containing the code defining the structure and style of the document
\usetikzlibrary{arrows.meta}
\usetikzlibrary{decorations}
\usetikzlibrary{decorations.markings}
\usetikzlibrary{decorations.pathreplacing}
\usetikzlibrary{shapes}
\usetikzlibrary{positioning}
\usetikzlibrary{backgrounds}

\begin{document}
\begin{tikzpicture}
	% Time axis:
	\node (time_origin) at (-2,4) {$t_{\Omega}$};
	\node (time_epoch_2) at (-2,2) {$t_{\text{upper}}$};
	\node (time_epoch_1) at (-2,1) {$t_{\text{lower}}$};
	% Nodes
	\node (origin_state) at (0,4) {$[1.0, 0.0]$};
	\draw (0, 2) node[minimum size=0pt,inner sep=0pt, outer sep=0pt] (epoch_boundary_2) {};
	\draw (0, 1) node[minimum size=0pt,inner sep=0pt, outer sep=0pt] (epoch_boundary_1) {};
	% Edges
	\draw[line width=1pt, dashed] (time_origin) -- (time_epoch_2) -- (time_epoch_1) -- (-2,-4);
	\draw [line width=3pt, white!50!black] (origin_state)  -- (epoch_boundary_2);
	\draw [line width=3pt, red!75!black] (epoch_boundary_2)  -- (epoch_boundary_1);
	\draw [line width=3pt, green!50!black] (epoch_boundary_1) -- (0, -4);
	% Labels
	\draw [decorate,decoration={brace,amplitude=10pt}] (2,4.3) -- (2,3.7) node [midway, xshift=1.25in, text width=2in, align=left] {Feature absent at origin};
	\draw [decorate,decoration={brace,amplitude=10pt}] (2,3.5) -- (2,2.1) node [midway, xshift=1.25in, text width=2in, align=left] {$r = 0$ (no change)};
	\draw [decorate,decoration={brace,amplitude=10pt}] (2,1.9) -- (2,1.1) node [midway, xshift=1.25in, text width=2in, align=left] {$\mathbf{Q} = \begin{bmatrix} -1 & 1 \\ 0 & 0 \end{bmatrix},\ r\text{ extremely small}$\\Feature arises in this date range\\(only once, with high probability)};
	\draw [decorate,decoration={brace,amplitude=10pt}] (2,0.9) -- (2,-3.9) node [midway, xshift=1.25in, text width=2in, align=left] {$\mathbf{Q} = \begin{bmatrix} 0 & 0\\ 1 & -1 \end{bmatrix},\ r\text{ moderate}$\\Feature cannot arise anymore;\\preserved in some branches\\and lost in others};
\end{tikzpicture}
\end{document}